\clearpage
% ====================================================================
% Section: Conclusions and Future Work
% ====================================================================
\chapter{Conclusions and Future Work}\label{cap:conclusion}

This dissertation has presented a kill-chain-aware adaptive defense framework for IoT networks that combines LSTM-based attack-stage modelling, supervised stage detection, and Deep Reinforcement Learning under a unified evaluation contract. The framework was developed and rigorously evaluated on the contemporary CICIoT2023 dataset across an audit-first, eight-stage development protocol. This concluding chapter summarises the contributions (Section~\ref{sec:summary}), discusses the empirical findings worth defending (Section~\ref{sec:findings}), surfaces the limitations of the approach (Section~\ref{sec:limitations}), and outlines a roadmap of future research directions (Section~\ref{sec:future_work}).

% --- Subsection 5.1 ---
\section{Summary of Contributions}\label{sec:summary}
The proliferation of IoT devices has created a vast and vulnerable attack surface, overwhelming conventional security paradigms that rely on static signatures and manual intervention. This dissertation has addressed this critical gap with a closed-loop security framework that synergises predictive analytics with autonomous, adaptive defense. The primary contributions are:
\begin{enumerate}
    \item \textbf{A kill-chain-aware adaptive defense architecture.} A two-stage framework integrating an LSTM-based red-team episode generator and a supervised stage detector with a five-action DRL defender (\textsc{observe}, \textsc{log}, \textsc{throttle}, \textsc{block}, \textsc{isolate}) trained against a stage-action proportionality reward. The architecture operates under an explicit \textbf{monotonic-attacker assumption} --- the red-team LSTM generates upper-triangular stage transitions, consistent with the IoTWarden threat model and the CICIoT2023 attack taxonomy --- and the defender observes only the 29-dimensional per-flow feature vector, never the true attack stage.

    \item \textbf{Empirical grounding on a contemporary IoT dataset, with an honest leakage-bug fix.} The framework is developed and validated on the \textbf{CICIoT2023 dataset}, projected onto a five-stage Cyber Kill Chain. A subtle splits-construction bug --- which had inadvertently leaked 70\,\% of every held-out out-of-distribution attack class back into the training pool --- was discovered during stage-detector training, fixed at commit \texttt{3cd2fb9}, and locked behind disjointness assertions before any benchmark numbers were collected. The fix and its forensics are presented honestly in Section~\ref{sec:dataset_results}.

    \item \textbf{A rigorous comparative DRL benchmark with a \LatencyRatio$\times$ latency advantage.} DQN, PPO, and A2C were trained over \NumSeeds{} seeds and evaluated on a held-out balanced test split ($n=300$ episodes per policy) against four reference policies (random, always-\textsc{observe}, always-\textsc{block}, RF-Acting) and an oracle Recommended-Action rule that observes the true attack stage. The best trained agent (A2C, $+\BestAgentReward$) captures \textbf{\OracleCapturePct\,\%} of the oracle ceiling ($+\OracleCeiling$) on the held-out test split, with non-overlapping bootstrap confidence intervals against every non-RL trivial baseline. Among deployable policies, trained A2C achieves \BestAgentLatency\,ms p50 inference latency versus RF-Acting's \RFLatency\,ms --- a \textbf{\LatencyRatio$\times$ latency advantage} at an approximately 12\,\% reward cost. The best-reward trained agent (A2C) exhibits benign FPR of 11.5\,\%; DQN offers the most FPR-conservative operating point at 6.1\,\%.


    \item \textbf{Reward mis-specification: discovered, diagnosed, and structurally fixed.} A 12-cell sparse one-at-a-time reward-component sweep identified that the reward-maximising policy under the naive reward contract ($\texttt{impact\_is\_terminal}=\mathit{True}$) diverges from the human-intended objective (``defend the \textsc{impact} step'') because de-escalation bonuses dominate the \textsc{impact}-step penalty. This is reward mis-specification in the formal sense~\cite{Amodei2016,Krakovna2020}. A \emph{structural} environment change --- setting $\texttt{impact\_is\_terminal}=\mathit{False}$ (the corrected primary contract) --- resolves the mis-specification at its source. An ablation probe (PPO, $n=30$ episodes) shows the structural fix lifts the mitigated-impact rate to $\FnineStructuralMitRate$ (ablation scale); at full benchmark scale ($n=300$, 10~seeds) the best trained agent (A2C) achieves a mitigated-impact rate of $0.317$. A six-point sweep over the defender de-escalation probability characterises the trained policy's response to attacker aggressiveness as monotone non-decreasing. Evaluation against four held-out OOD attack classes activates the pre-registered finding \texttt{D7.9.1}: the trained policy (best: PPO $+1355.2$) is \emph{robust to} (within seed-noise of its in-distribution mean), but not \emph{better at}, the hardest OOD class on which the supervised stage detector is structurally blind (RF-Acting $+1680.0$, $\Delta=-324.8$).

    \item \textbf{An audit-first reproducibility protocol.} Every figure ships a \texttt{manifest.json} that pins the SHA-256 hashes of every input artefact and the producing Git commit. A smoke-reproducibility harness (\texttt{python -m scripts.reproducibility\_smoke}) verifies the full hash chain end-to-end on a fresh checkout in approximately five seconds, returning a verdict bucketed into \textsc{ok} / \textsc{fail} / \textsc{known-divergence} / \textsc{skip}. The complete protocol is documented in Appendix~A.
\end{enumerate}

% --- Subsection 5.2 ---
\section{Principal Empirical Findings}\label{sec:findings}
Three findings are pre-registered, principled results of the approach rather than ad-hoc failures, and are arguably the strongest contributions of this dissertation:

\begin{itemize}
    \item \textbf{Finding 1 (Blue-Team Training stage, gate G5.4 \textsc{pass-with-finding}): Reward Mis-specification, Discovered and Structurally Fixed.} A reward-decomposition trace explains the gap between high episodic reward and low mitigated-impact rate as principled reward mis-specification~\cite{Amodei2016,Krakovna2020} under the naive contract ($\texttt{impact\_is\_terminal}=\mathit{True}$): de-escalation bonuses earned during the kill chain dominate the expected \textsc{impact}-step loss, so the reward-maximising policy --- ``rack up de-escalations and accept the \textsc{impact} loss'' --- is \emph{correct} under the proxy objective but \emph{wrong} under the intended objective. The finding is not a regression: it is the input that motivates the reward-component ablation, which identifies and applies the structural fix ($\texttt{impact\_is\_terminal}=\mathit{False}$, the corrected primary contract). An ablation probe (PPO, $n=30$) shows the structural fix lifts the mitigated-impact rate to \FnineStructuralMitRate; at full benchmark scale ($n=300$, 10~seeds) the best trained agent (A2C) achieves $0.317$. The FPR trade-off remains: A2C exhibits 11.5\,\% benign FPR, underscoring that a Lagrangian FPR penalty or hierarchical policy is needed to reach operational deployment thresholds.

    \item \textbf{Finding 2 (Held-Out Benchmark stage, audit AF2 reframe of gate G6.2): \OracleCapturePct\,\% of the Oracle Ceiling.} The best trained DRL agent (A2C, $+\BestAgentReward$) captures \OracleCapturePct\,\% of the oracle ceiling ($+\OracleCeiling$ on \texttt{test\_balanced}) without observing the true attack stage; the oracle is reported as a measurement instrument rather than a competing baseline because it reads the true attack stage directly from the simulator and cannot be deployed. The AF2 reframe is preserved in the JSON scoreboard record (the original ``RL $>$ recommended-action rule'' threshold reads \texttt{passes: false} permanently), so it is not goalpost-moving --- it is an interpretation correction documented in the results prose. All three RL agents are statistically indistinguishable from each other on mean reward (n=10 seeds, overlapping CIs); the main deployable trade-off is between RF-Acting (higher reward, \RFLatency\,ms latency) and trained RL (lower reward, \BestAgentLatency\,ms latency, \LatencyRatio$\times$ faster).

    \item \textbf{Finding 3 (Ablation \& Robustness stage, gates G7.2 \textsc{pass-with-finding} and G7.9 \textsc{fail-with-finding}): Structural Fix and OOD Robustness Limitation.} A 12-cell reward-coefficient sweep characterises the limit of one-at-a-time coefficient shaping: 11 of 12 cells stay within $\pm 150$ reward of the centre baseline. The structural fix ($\texttt{impact\_is\_terminal}=\mathit{False}$) achieves an ablation-probe mean reward of $+\FnineStructuralReward$ (PPO, $n=30$) and mitigated-impact rate \FnineStructuralMitRate. Out-of-distribution evaluation against the hardest held-out class \texttt{VulnerabilityScan} (supervised-detector recall $0.001$) activates the pre-registered \texttt{D7.9.1}: the best trained RL policy (PPO $+1355.2$, CI [1312, 1393]) does not collapse on this class but does not exceed RF-Acting ($+1680.0$, CI [1641, 1720]; $\Delta=-324.8$). Closing the OOD gap requires either explicit attack-class curriculum or train-time OOD-augmented data, reframed as a limitation (Limitation~3 below and Direction~3 in Section~\ref{sec:future_work}).
\end{itemize}

These three findings, together with the audit-first reproducibility protocol that makes them verifiable, constitute the empirical backbone of the dissertation.

% --- Subsection 5.3 ---
\section{Limitations}\label{sec:limitations}
The work has five principal limitations, surfaced honestly in Section~\ref{sec:discussion} and recapitulated here.

\begin{enumerate}
    \item \textbf{Compromise rate $\equiv 1.0$ on \texttt{test\_balanced}.} Under \emph{both} contracts ($\texttt{impact\_is\_terminal}=\mathit{True}$ and $=\mathit{False}$), the upper-triangular red-team LSTM eventually advances every episode to \textsc{impact}; defender-driven de-escalation only resets the chain to \textsc{benign}. The meaningful security KPI throughout this dissertation is therefore \texttt{mitigated\_impact\_rate}, not \texttt{compromise\_rate}. The $\texttt{impact\_is\_terminal}=\mathit{False}$ contract enables reactive mitigation at the \textsc{impact} step but does not eliminate the episode-level compromise. The DRL policies trained here constitute \emph{reactive mitigation} agents, not \emph{pre-emption} agents.

    \item \textbf{MTTC values bounded by the lifecycle-floor \textsc{impact} clamp.} The clamp downgrades pre-step-20 \textsc{impact} transitions to \textsc{maneuver}, so mean MTTC is biased toward $\texttt{min\_episode\_length} = 20$. Mean MTTC numbers should not be interpreted as ``unconstrained natural'' MTTC.

    \item \textbf{Trained RL does not beat RF-Acting on \texttt{test\_balanced} or on the hardest OOD class.} The deployable comparison is a latency--reward trade-off frontier: RF-Acting achieves higher reward ($+\RFReward$) at \LatencyRatio$\times$ higher inference latency (\RFLatency\,ms vs.\ \BestAgentLatency\,ms). The dissertation makes the trade-off explicit and frames it honestly; it does not claim RL dominates RF-Acting in absolute reward. On the hardest OOD class (\texttt{VulnerabilityScan}), RF-Acting widens the gap to $\Delta=-324.8$.

    \item \textbf{Elevated benign false-positive rate.} All trained agents exhibit benign-FPR between 6.1\,\% (DQN) and 11.5\,\% (A2C), substantially above the 1\,\% operational threshold. The best-reward agent (A2C) carries the \emph{highest} FPR, making A2C the highest-reward but least FPR-conservative choice. Addressing this requires either a hard FPR constraint in the reward function (the Lagrangian ablation explored in this dissertation is a first step) or a hierarchical policy that separates stage detection from action selection. This limitation should be a primary consideration for any real deployment decision.

    \item \textbf{Upper-triangular red-team LSTM (monotonic-attacker assumption).} The red-team LSTM does not allow stage retreats, limiting realism relative to a real adversary that may pivot, retreat, or interleave kill-chain stages. This is a deliberate, explicitly-bounded simplification consistent with the IoTWarden threat model; non-monotonic attacker modelling is deferred to future work.
\end{enumerate}

% --- Subsection 5.4 ---
\section{Future Work and Research Directions}\label{sec:future_work}
The audit-first protocol that produced this dissertation explicitly delineated work scoped \emph{into} this dissertation from work scoped \emph{out of} it. Six follow-up directions are surfaced here in order of mentor-recommended priority.

\subsection{Tightly Scoped Extensions of Empirical Findings}

\paragraph{Direction 1 --- Non-monotonic attacker model.} The current red-team LSTM is upper-triangular by construction: the kill chain advances but never retreats. Relaxing this to a non-monotonic LSTM (or Transformer-based) attacker that allows stage pivots and retreats would produce more realistic episode trajectories and test the defender's robustness to adversarial pivoting. This is the most direct extension of the monotonic-attacker limitation, as the training environment's reward function already supports de-escalation; only the transition generator needs to change. It also directly addresses the threat-model boundary identified in Section~\ref{ssec:monotonic_attacker}.

\paragraph{Direction 2 --- Edge-hardware deployment and latency budget.} The \LatencyRatio$\times$ latency advantage of trained A2C over RF-Acting (\BestAgentLatency\,ms vs.\ \RFLatency\,ms p50) establishes trained RL as the candidate policy for resource-constrained edge-hardware deployments. A natural follow-up is to quantify the latency and accuracy trade-off on actual IoT gateway hardware (e.g.\ ARM Cortex-A-class processors) under real packet rates, and to apply model-compression techniques (quantisation, pruning) to the trained policy network to fit within the tighter computational budgets of embedded defenders. The policy's $\sim$4K-parameter footprint ($\sim$20\,KB) already makes it competitive with embedded ML models.

\paragraph{Direction 3 --- Train-time OOD-class augmentation.} The complement of the F15 OOD-robustness evaluation: instead of evaluating against held-out classes, \emph{include} a simulacrum of them in training (synthetic feature blending, domain randomisation, or a small adversarial-augmentation generator) and check whether the resulting policy then exceeds RF-Acting on \texttt{VulnerabilityScan}. This directly tests whether the ``robust to, not better at, the OOD class'' framing is a structural limit or a curable training-data limitation (the \texttt{D7.9.1} finding-activation left this question open).

\subsection{Broader Research Directions}

\paragraph{Direction 4 --- Multi-agent reinforcement learning and self-play adversary.} A real IoT deployment includes many overlapping defenders (per-segment defenders coordinating with a central Security Operations Centre). This can be modelled as a cooperative multi-agent RL (MARL) system where specialised agents defend different network segments and learn to coordinate. Extending to a self-play adversary --- where the red-team LSTM is jointly trained with the defender to produce adversarial stage trajectories --- would produce a harder, more dynamic training distribution that could close the remaining gap to the oracle ceiling without a structural reward-function change. Both extensions inherit the kill-chain decision frame and the audit-first reproducibility protocol from this dissertation.

\paragraph{Direction 5 --- Federated learning for multi-site deployment.} A single-site RL policy trained on one IoT deployment may not generalise to a different network topology or device mix. Federated extensions would allow defender policies to be trained across diverse IoT deployments without centralising sensitive traffic data, producing a population-level policy that is more robust to site-specific distribution shift. This direction is complementary to Direction 4 and is directly motivated by the OOD evaluation's finding that per-site feature novelty limits generalisation.

\paragraph{Direction 6 --- Reward-function redesign with explicit FPR constraints.} The elevated benign-FPR (6.1--11.5\,\%) is a direct consequence of the unconstrained de-escalation-bonus farming mechanism. A principled redesign would introduce a hard constraint on the benign-FPR (e.g.\ as a Lagrange multiplier in a constrained MDP) or a hierarchical policy that separates the stage-detection decision from the action-selection decision. This direction is the most directly actionable follow-up to Limitation 4 and would produce a policy that is safe to deploy in environments where false positives have high operational cost.

\subsection{Note on the Environment-Perturbation Robustness Stage}
The audit-first protocol that structured this dissertation defined eight development stages (stages 1--7 delivered; stage 8 reserved for environment-perturbation robustness work including noise-robustness (F13) and OOD-augmentation (F14)). The environment-perturbation robustness stage was \emph{skipped} as a dissertation deliverable: the F9 structural change resolved the reward mis-specification using only the $\texttt{impact\_is\_terminal}$ flag, making a full perturbation sweep less informative for the dissertation's headline claim than a precise articulation of the existing ablation-stage findings. The F13 and F14 directions are reframed above as Directions 3 and 2 (post-thesis future work). This decision is part of the audit-trail anchored at the four committed \texttt{G[N]\_scoreboard.json} files.
